\section{相关工作}
\label{sec:background}

\subsection{直接偏好优化及其变体}

直接偏好优化 \cite{rafailov2023dpo} 通过将奖励函数隐式地重参数化为策略本身,重新表述了 RLHF,从而免去了独立奖励模型的需求。从 Bradley-Terry 偏好模型推导出的闭式解使训练得以稳定,但也在策略的似然景观与隐式奖励曲面之间引入了一种根本性的耦合。恒等偏好优化(IPO)\cite{azar2024ipo} 通过优化一个避免对数比饱和问题的平方损失,来应对 DPO 对确定性偏好的过拟合;Kahneman-Tversky 优化(KTO)\cite{ethayarajh2024kto} 则借助前景理论中对单个输出的价值函数,彻底免去了对成对偏好的需求。SimPO \cite{meng2024simpo} 通过使用序列级平均对数概率作为隐式奖励来去除对参考模型的依赖;ORPO \cite{hong2024orpo} 则在单一阶段中统一了指令微调与偏好对齐。尽管这些算法各有改进,所有变体都共享一个共同的结构性质:它们在偏好对上优化一个对比式或基于间隔的目标,从而产生强烈的梯度压力,将质量推向一小撮高奖励响应。这种朝向模式集中的共有归纳偏置——尽管作为训练稳定性问题被研究过——尚未被作为对抗利用的载体加以审视。

\subsection{偏好学习中的模式坍缩}

经偏好优化的模型坍缩到狭窄输出分布上的倾向,作为一个\emph{质量}问题已有充分记载。Razin 等人 \cite{razin2024vanishing} 提供了一项理论分析,表明强化学习微调会随策略偏离参考模型而出现梯度消失,使训练在退化模式周围停滞。Pal 等人 \cite{pal2024smaug} 在 DPO 中识别出一种具体的失效模式:当被选响应的似然低于被拒响应时会产生反向梯度,并提出 DPO-Positive 以过滤此类病态对。Kirk 等人 \cite{kirk2024diversity} 经验性地证明 RLHF 会跨领域系统性地降低输出多样性,表明对齐过程以一种跨提示策略持续存在的方式将模型行为同质化。关键在于,这整条研究脉络都仅从模型能力退化的视角来框定模式坍缩:多样性降低损害创造性任务,退化模式降低指令遵循质量,梯度病态减慢收敛。此前没有任何工作认识到,模式坍缩所引入的\emph{可预测性}——即坍缩模型将概率质量集中到一个狭窄、可刻画的输出集合上这一事实——构成了一项与安全相关的性质,攻击者可以蓄意诱导并随后加以利用。

\subsection{对大模型对齐的投毒攻击}

越来越多的工作表明偏好学习流程易受数据投毒。Pathmanathan 等人 \cite{pathmanathan2025dpo_poison} 表明,仅破坏 0.5\% 的 DPO 训练对就足以植入绕过安全过滤器的定向不安全行为,在有害查询上达到超过 20\% 的攻击成功率。Wang 等人 \cite{wang2024rlhfpoison} 攻击 RLHF 流程中的奖励模型,证明奖励投毒会通过策略优化传播,产生可靠的不安全输出。影子对齐 \cite{yang2023shadow} 仅用 100 个对抗样本,借助微调的灾难性遗忘动力学实现安全性退化;Souly 等人 \cite{souly2025nearconstant} 则确立,无论数据集规模如何,近乎恒定数量的投毒样本即已足够。最近,颠覆性对齐注入(SAI)\cite{mamun2025sai} 证明对抗偏好对可通过利用 DPO 损失基于间隔的结构来覆盖安全训练。尽管这些攻击令人信服地证明了 DPO 训练易受对抗数据影响,它们无一例外地将投毒机制视为一个直接的输入-输出映射问题:注入有害对,产出有害输出。这些工作中没有任何一个识别或利用偏好学习的\emph{模式坍缩动力学}作为放大机制——它们投毒的是模型关于何为安全的知识,而非其分布几何。

\subsection{越狱攻击与防御}

推理期的越狱攻击代表了一种与训练期投毒互补的威胁模型。贪婪坐标梯度(GCG)\cite{zou2023gcg} 通过离散 token 优化构造可在模型间迁移的对抗后缀;AutoDAN \cite{liu2024autodan} 使用分层遗传算法自动生成越狱,产出语义上有意义的攻击字符串。PAIR \cite{chao2024pair} 与带剪枝的攻击树(TAP)\cite{mehrotra2024tap} 利用攻击者大模型通过多轮交互迭代精化越狱提示,在无需梯度访问的情况下取得很高的成功率。在防御一侧,SmoothLLM \cite{robey2024smoothllm} 施加随机的字符级扰动,通过输出的不一致性来检测对抗后缀;RA-LLM \cite{cao2024rallm} 通过基于检索的接地来增强鲁棒性;困惑度过滤 \cite{jain2023perplexity} 则拒绝困惑度异常高的输入。这些方法完全在推理期运作,利用模型的安全训练与其在分布偏移下实际行为之间的差距。然而,它们并不与训练动力学交互:一个通过对抗训练抵御 GCG 攻击的模型,对偏好坍缩利用同样脆弱,因为我们识别出的脆弱性存在于优化景观中,而非输入空间中。

\subsection{本工作的定位}

上述三类文献处理了对齐脆弱性中相邻却彼此脱节的方面。模式坍缩研究认识到偏好优化会集中输出分布,但仅将其视为一个可通过正则化处理的质量退化问题。投毒研究证明对抗训练数据可以损害安全性,但它通过直接的行为操纵进行运作,既未利用——甚至未承认——坍缩在输出分布中所诱导的几何结构。越狱研究利用推理期的行为差距,却无法访问或修改训练动力学。我们的工作识别出一个此前未被认识到的联系:\emph{模式坍缩不仅是一个性能缺陷,更是一项可被系统性利用的安全漏洞}。我们证明:攻击者可以构造蓄意将坍缩加速引向攻击者所选模式的投毒数据;由此产生的分布集中使模型的不安全行为变得可预测、可触发;并且标准的坍缩缓解技术(KL 惩罚、多样性正则化)是不充分的,因为它们是为保持质量而非对抗鲁棒性设计的。这一重新框定——从模式坍缩作为不便,到模式坍缩作为攻击面——在优化理论与 AI 安全的交叉处开启了一个新的研究方向。
